% ============================================================================
% Fichier cible : ARIEL_LOGMINER_MEMOIRE_FINAL/chapters/chapitre2_etat_art.tex
% Passage à remplacer : de « \section{Systèmes multi-agents} » (inclus) à
%   « \section{Routage et spécialisation des traitements} » (exclu), donc les
%   sections actuelles « Systèmes multi-agents » et « Bus de messages... ».
% Agent : Agent 3 — État de l'art et positionnement, chantier 2026-11-05.
% Sources : Wooldridge ; Smith ; FIPA ; documentations officielles Celery, Kafka,
%   Redis Streams et Temporal.
% ============================================================================

\section{Agents, files de tâches et moteurs de workflow}

\subsection{Ce qui distingue théoriquement un agent d'un worker}

Un système multi-agent décompose un problème entre des entités qui possèdent un état,
des capacités et une politique de décision. Wooldridge associe classiquement l'agent à
l'autonomie, la réactivité, la proactivité et la capacité sociale
\cite{wooldridge2009introduction}. Ces propriétés excèdent la simple exécution d'une
fonction reçue dans une file. Elles ne sont cependant pas acquises par le vocabulaire
employé : un worker muni d'une table de capacités et d'un score fixe reste
fonctionnellement proche d'un ordonnanceur déterministe si ses décisions ne diffèrent pas
de celles de cette table.

Le protocole Contract Net de Smith organise une allocation négociée : un gestionnaire
annonce une tâche, les contractants potentiels soumettent une offre, puis une offre est
retenue \cite{smith1980contract}. La spécification FIPA précise les actes
\texttt{cfp}, \texttt{propose}, \texttt{refuse}, \texttt{accept-proposal},
\texttt{reject-proposal} et les notifications finales
\cite{fipa2002contractnet,fipa2002aclmessage,fipa2002cal}. Le protocole d'Ariel Logminer
en reprend la structure sous une nomenclature propre, avec des extensions
\texttt{ACCEPT} et \texttt{FEEDBACK} ; il revendique une parenté, non une conformité
FIPA.

\subsection{Le concurrent réel : l'affectation déterministe}

Pour une chaîne de journaux, Contract Net doit être comparé non seulement à un monolithe,
mais à des mécanismes d'orchestration plus simples. Celery associe des noms de tâches à
des fonctions et permet aux workers de consommer des files ; Kafka répartit les
enregistrements entre consommateurs et documente explicitement les garanties de
livraison ; Redis Streams maintient des groupes de consommateurs et une liste des
messages livrés mais non acquittés ; Temporal persiste l'historique d'un workflow et
place ses unités de travail dans des files interrogées par des workers
\cite{celery2026tasks,kafka2026semantics,redis2026streams,temporal2026tasks}.
Dans chacun de ces cadres, une table de routage, des files nommées ou l'inscription de
workers par capacité peuvent affecter une tâche sans appel d'offres.

\begin{table}[H]
\centering
\caption{Contract Net et alternatives d'orchestration}
\label{tab:concurrents-orchestration}
\scriptsize
\begin{tabularx}{\textwidth}{p{2.3cm}p{3cm}p{3.2cm}X}
\toprule
\textbf{Mécanisme} & \textbf{Décision d'affectation} & \textbf{État et reprise} &
\textbf{Question comparative} \\
\midrule
Table statique / files nommées & Règle source--pipeline ou type--file fixée à l'avance. &
Dépendent du broker et du magasin applicatif. &
La négociation modifie-t-elle l'affectation ou seulement son protocole ? \\
Groupe de consommateurs & Le broker distribue les messages entre workers d'un groupe. &
Messages en vol, acquittements et réattribution selon le broker. &
Une utilité locale améliore-t-elle débit, équité ou compatibilité ? \\
Moteur de workflow & Graphe ou code de workflow durable, tâches placées dans des files. &
Historique persistant, reprises et tentatives explicites. &
Le modèle par objectifs apporte-t-il une propriété absente du workflow ? \\
Contract Net & Appel à propositions, offres, puis attribution. &
À construire au-dessus du transport et des magasins. &
Le coût des offres est-il compensé par de meilleures décisions ? \\
\bottomrule
\end{tabularx}
\end{table}

Une justification scientifique de l'architecture à agents exige donc un comparateur
contrôlé : mêmes tâches, mêmes workers, même transport, mêmes mécanismes de reprise et
seule politique d'affectation différente. Sans cette ablation, la modularité du code et
la présence de messages Contract Net démontrent une réalisation, pas la supériorité de la
négociation.

\section{Sémantique de livraison, reprise et idempotence}

Les garanties de transport doivent être séparées des garanties métier. Kafka distingue
les sémantiques \textit{au plus une fois}, \textit{au moins une fois} et
\textit{exactement une fois}, et limite sa garantie exactement une fois au périmètre de
transactions et de traitements compatibles \cite{kafka2026semantics}. Dans Redis
Streams, un groupe conserve les messages non acquittés dans une liste d'attente ; un
message peut être relu ou réattribué après une panne, donc traité plusieurs fois
\cite{redis2026streams}. Celery formule la même contrainte : lorsque l'acquittement a lieu
après l'exécution, une panne peut provoquer une nouvelle exécution et la tâche doit être
idempotente \cite{celery2026tasks}.

L'idempotence est ainsi une propriété de l'opération ou de son magasin, non une propriété
spécifique des agents. Elle consiste à rendre plusieurs tentatives équivalentes à une
seule du point de vue de l'effet observable, par exemple au moyen d'une clé stable et
d'une écriture conditionnelle. Le broker peut préserver et réattribuer une tâche ; il ne
peut pas, à lui seul, annuler un effet externe déjà produit. Inversement, un magasin
d'idempotence peut protéger un monolithe, des workers centralisés ou des agents.

Une preuve de reprise doit donc documenter séparément : la livraison ou la réattribution,
le nombre de tentatives, l'effet métier, l'acquittement et la persistance du résultat.
Elle ne vaut que pour les pannes effectivement injectées. Démontrer la reprise d'un
worker ne démontre ni la tolérance à une panne du broker, ni le fonctionnement pendant
une partition réseau, ni une sémantique exactement une fois de bout en bout.

